\documentclass[10pt,a4paper]{amsart}
\usepackage[margin=19mm]{geometry}
\usepackage{amsmath,amssymb,mathtools}
\usepackage[scr=rsfs,cal=euler]{mathalfa}
\usepackage{xfrac}
\usepackage[unicode,hidelinks,bookmarks=false]{hyperref}
\newcommand{\ie}{i.e.\ }
\newcommand{\cf}{cf.\ }
\newcommand{\resp}{respectively\ }
\newcommand{\Subsection}{\subsection}
\providecommand{\nobreakdash}{\nobreak\hbox{-}}
\setlength{\emergencystretch}{2em}
\multlinegap0pt
\usepackage{bm}
\usepackage{bbm}
\usepackage{dsfont}
\usepackage{xspace}
\usepackage{cancel}
\usepackage{mathtools}
\usepackage{amsthm}
\makeatletter
\@ifundefined{lemma}{\newtheorem{lemma}{Lemma}}{}
\@ifundefined{remark}{\newtheorem{remark}{Remark}}{}
\makeatother
\title[Scaling Optimized Spectral Approximations on Unbounded Domai]{Scaling Optimized Spectral Approximations on Unbounded Domains: The Generalized Hermite and Laguerre Methods}
\author{Hao Hu, Haijun Yu}
\date{Source: arXiv:2602.03083v1 (CC BY 4.0)}
\begin{document}
\maketitle

\noindent\emph{Source and licence.} Scaling Optimized Spectral Approximations on Unbounded Domains: The Generalized Hermite and Laguerre Methods. Version arXiv:2602.03083v1 (CC BY 4.0), distributed under the Creative Commons Attribution 4.0 International licence, as recorded in the arXiv licence metadata for this exact version and shown on the abstract page https://arxiv.org/abs/2602.03083v1. Full rights notice: licenses/arxiv-2602.03083-CC-BY-4.0.txt.\par\smallskip

\noindent\textit{Editorial scope.} Counted unit: the Remark whose source label is \texttt{None} in the article (source0.tex lines 467--481, source label \texttt{None}), delivered together with the article's own complete original proof of it (source0.tex lines 483--759, one \texttt{proof} environment of 277 lines). No mathematics is written, repaired or re-derived: statement and proof are the article's own text line for line. Required context retained uncounted: lem:transform (lines 338--347); eq:bigthm (lines 456--463). Every definition, display and label of those lines is retained unchanged. Omitted from this package and not redistributed: 1--337, 348--455, 464--466, 482--482, 760--2167. Nothing inside a retained excerpt is omitted or altered: every retained body is delivered exactly as the article's lines are written, so no omission receipt of any kind accompanies this package. Citations: the retained excerpts carry no citation command at all, so no bibliography entry of the article is transcribed and no reference is redistributed. Reference adaptation: none was needed and none was made; every reference inside the delivered unit resolves inside it, so no unresolved reference or citation is produced. Printed slips kept literal: no printed word, symbol, hypothesis or number of the counted statement or of its proof was altered. Layout and class: the article's own class is \texttt{sn-jnl} with packages including graphicx, multirow, amsmath, amssymb, amsfonts, amsthm, mathrsfs, appendix, xcolor, textcomp, manyfoot, booktabs, algorithm, algorithmicx; the wrapper is the lane's pinned \texttt{amsart} editorial wrapper at 10pt on A4, which restates the theorem-like environments the retained text needs and adds the article's own macro definitions that the retained text needs (none), with extra wrapper packages (bm, bbm, dsfont, xspace, cancel, mathtools). The delivered numbering is the unit's own. Assets: no retained span contains a figure environment, a graphics-inclusion command, a PDF-page inclusion, a table or a TikZ picture; no image file is copied into this package, the package declares no asset dependency and no image file is redistributed with it. Licence: version arXiv:2602.03083v1 of this article is distributed under the Creative Commons Attribution 4.0 International licence, as recorded in the arXiv licence metadata for this exact version and shown on the abstract page https://arxiv.org/abs/2602.03083v1. A full rights notice is delivered at licenses/arxiv-2602.03083-CC-BY-4.0.txt. Characters: every retained excerpt is pure ASCII, so the delivered engine, XeLaTeX with its default Latin Modern text font, reports no missing character.
\begin{lemma}\label{lem:transform}
  \begin{equation}\label{eq:0823-2}
    \begin{aligned}
    & \sqrt{2}\mathcal{F}\left[e^{-u^2}|u|^{2 \mu} H_{2 m}^\mu(u)\right](-2x)= \\
    & \qquad\qquad(-1)^m 2^{2 m} \frac{\Gamma\left(m+\mu+\frac{1}{2}\right)}{\Gamma\left(m+\frac{1}{2}\right)} x^{2 m} e^{-x^2} \Phi\left(-\mu, m+\frac{1}{2} ; x^2\right), \\
    & \sqrt{2}\mathcal{F}\left[e^{-u^2}|u|^{2 \mu} H_{2 m+1}^\mu(u)\right](-2x)= \\
    & \qquad \qquad(-1)^m 2^{2 m+1} \frac{\Gamma(m+\mu+3 / 2)}{\Gamma(m+3 / 2)} i\, x^{2 m+1} e^{-x^2} \Phi\left(-\mu, m+3 / 2 ; x^2\right) .
    \end{aligned}
  \end{equation}
\end{lemma}

  \begin{equation}\label{eq:bigthm}
    \begin{aligned}
    \left\|u-\widehat{\Pi}_N^\beta u\right\|_\omega & 
    \lesssim_\mu \bigl\|u \cdot \mathbb{I}_{\left\{|x|>M\right\}}\bigr\|_\omega 
     +B^{-\mu+\frac{1}{2}} \bigl\| \mathcal{F}[u]\left(B \xi\right) \bigr\|_{H^{\mu}\left(\mathbb{R}\backslash [-1,1]\right)} \\
    &{} \quad +e^{-\frac{N}{24}}\left(\|u\|_\omega+\|u\| \cdot M^\mu\right),
    \end{aligned}
  \end{equation}

\begin{remark}
  This theorem can be regarded as an analogue of the Nyquist-Shannon sampling theorem: 
  When using the first $N+1$ scaled Hermite functions $\widehat{H}_n^{(\mu)}(\beta x)$, the capturable bandwidth ranges in the spatial and frequency domains are $M = \mathcal{O}\bigl(\sqrt{N}/\beta\bigr)$ and $B = \mathcal{O}\bigl(\sqrt{N}\beta\bigr)$, respectively. 
  The approximation error arises from three sources: 
  \begin{enumerate}
    \item $\left\|u \cdot \mathbb{I}_{\left\{|x|>M\right\}}\right\|_\omega$:
    The insufficient decay of components beyond the spatial bandwidth.
    \item $B^{-\mu+\frac{1}{2}} \left\| \mathcal{F}[u]\left(B \xi\right) \right\|_{H^{\mu}\left(\mathbb{R}\backslash [-1,1]\right)}$:
    the high-frequency components beyond the frequency bandwidth.
    \item $e^{-\frac{N}{24}}\left(\|u\|_\omega+\|u\| \cdot M^\mu\right)$:
    The exponentially decaying error that is attributed to the approximation 
    within the truncated bandwidth range.
  \end{enumerate}
  The following proof strategy is built upon this intuitive understanding.
\end{remark}

\begin{proof}
  Let $v(x) = u(x/\beta)$, then
  \begin{equation}\label{eq:0505-2}
    \left\|u-\widehat{\Pi}_N^\beta u\right\|_\omega=\beta^{-\mu-\frac{1}{2}}\left\|v-\widehat{\Pi}_N^1 v\right\|_\omega.
  \end{equation}
  Similarly, after replacing $u$ with $v$ and $\beta$ with $1$, 
  each term on the right-hand side of \eqref{eq:bigthm} 
  will also acquire the same additional factor $\beta^{-\mu-\frac{1}{2}}$,
  hence it is sufficient to prove \eqref{eq:bigthm} for $\beta = 1$.
  In this case, we have $M = B = \frac{\sqrt{N}}{2\sqrt{3}}$.

  Let us define the spatial truncation function as
  \begin{equation}\label{eq:0505-3}
    T_M^s=\mathbb{I}_{[-2 M, 2 M]} * \frac{1}{\sqrt{2 \pi}} e^{-\frac{1}{2} x^2}.
  \end{equation}
  Then define
  \begin{equation}\label{eq:0505-3.5}
    u_M(x) = u(x) \cdot T_M^s(x).
  \end{equation}
  Let
  \begin{equation}\label{eq:0505-4}
    \alpha(x)=\left\{\begin{array}{cc}
    e^{\frac{1}{x^2-1}} & |x|<1 \\
    0 & |x| \geqslant 1,
    \end{array}\right.
  \end{equation}
  and
  \begin{equation}\label{eq:0505-5}
    \gamma(x)=\left(\int_\mathbb{R} \alpha(x) d x\right)^{-1} \alpha(x) ; \quad \gamma_{\varepsilon}(x)=\varepsilon^{-1} \gamma\left(\frac{x}{\varepsilon}\right).
  \end{equation}
  Then, define the frequency truncation function
  \begin{equation}\label{eq:0505-6}
    T_B^f=\mathbb{I}_{\left[-\frac{3}{2} B, \frac{3}{2} B\right]} * \gamma_{\frac{B}{2}}.
  \end{equation}
  Further define
  \begin{equation}\label{eq:0505-7}
    u^B(x)=\frac{1}{\sqrt{2 \pi}} \int \mathcal{F}[u](\xi) \cdot T_B^f(\xi) \cdot e^{i \xi x} d \xi,  
  \end{equation}
  and
  \begin{equation}
  u_M^B=u^B \cdot T_M^S.
  \end{equation}
  With the above preparation, we can perform the following estimation. Since $\beta = 1$, we
  will ignore the superscript $\beta$ in $\widehat{\Pi}_N^\beta$ for simplicity.
  \begin{equation}\label{eq:0505-9}
    \begin{aligned}
    \left\|u-\widehat{\Pi}_N u\right\|_\omega & \leqslant\left\|u-u_M^B\right\|_\omega+\left\|u_M^B-\widehat{\Pi}_N u_M^B\right\|_\omega+\left\|\widehat{\Pi}_N\left(u_M^B-u\right)\right\|_\omega \\
    & \lesssim\left\|u-u_M^B\right\|_\omega+\left\|u_M^B-\widehat{\Pi}_N u_M^B\right\|_\omega \\
    & \lesssim\left\|u-u_M\right\|_\omega+\left\|u_M-u_M^B\right\|_\omega+\left\|u_M^B-\widehat{\Pi}_N u_M^B\right\|_\omega \\
    & \triangleq E_s+E_f+E_H.
    \end{aligned}
  \end{equation}
  Notice that
  \begin{equation}\label{eq:0505-10}
    \begin{aligned}
    E_s^2 & =\int\left(u-u_M\right)^2 \cdot \omega \\
    & =\int_{|x| \leqslant M}\left(u-u_M\right)^2 \cdot \omega+\int_{|x|>M}\left(u-u_M\right)^2 \cdot \omega.
    \end{aligned}
  \end{equation}
  For the first term,
  \begin{equation}\label{eq:0505-11}
    \int_{|x| \leqslant M}(u-u_M)^2 \cdot \omega=\int_{|x| \leqslant M} u^2\left(1-T_M^s\right)^2 \cdot \omega.
  \end{equation}
  Since for $|x| \leqslant M$
  \begin{equation}\label{eq:0505-12}
    \begin{aligned}
    \left|1-T_M^s\right| & =\left(\int_{2 M-x}^{+\infty}+\int_{-\infty}^{-2 M-x}\right) \frac{1}{\sqrt{2 \pi}} e^{-\frac{1}{2} t^2} d t \\
    & \lesssim \int_M^{+\infty} \frac{1}{\sqrt{2 \pi}} e^{-\frac{1}{2} t^2} d t \\
    & \lesssim e^{-\frac{1}{2} M^2},
    \end{aligned}
  \end{equation}
  putting \eqref{eq:0505-12} into \eqref{eq:0505-11} yields
  \begin{equation}\label{eq:0505-13}
    \begin{aligned}
    \int_{|x| \leqslant M}\left(u-u_M\right)^2 \cdot \omega & \lesssim e^{-M^2} \int u^2 \cdot \omega =e^{-M^2}\|u\|_\omega^2.
    \end{aligned}
  \end{equation}
  For the second term of \eqref{eq:0505-10} we have
  \begin{equation}\label{eq:0505-14}
    \int_{|x|>M}\left(u-u_M\right)^2 \cdot \omega=\int_{|x|>M} u^2\left(1-T_M^s\right)^2 \cdot \omega.
  \end{equation}
  Notice that $0 \leqslant T_M^s \leqslant 1$, we have
  \begin{equation}\label{eq:0505-15}
    \begin{aligned}
    \int_{|x|>M}\left(u-u_M\right)^2 \cdot \omega & \leqslant \int_{|x|>M} u^2 \cdot \omega 
     =\left\|u \cdot \mathbb{I}_{\{|x|>M\}}\right\|_\omega^2.
    \end{aligned}
  \end{equation}
  Putting \eqref{eq:0505-13}, \eqref{eq:0505-15} back into \eqref{eq:0505-10}
  yields
  \begin{equation}\label{eq:0505-16}
    \begin{aligned}
    E_s & \lesssim e^{-\frac{1}{2} M^2}\|u\|_\omega+\|u \cdot \mathbb{I}_{\{|x|>M\}}\|_\omega \\
    & =e^{-\frac{N}{24}}\|u\|_\omega+\left\|u \cdot \mathbb{I}_{\left\{|x|>M\right\}}\right\|_\omega.
    \end{aligned}
  \end{equation}
  For $E_f$ defined in \eqref{eq:0505-9} we have
  \begin{equation}\label{eq:0505-17}
    \begin{aligned}
    E_f^2 & =\int\left|u_M-u_M^B\right|^2 \cdot \omega \\
    & =\int\left|u-u^B\right|^2 \cdot\left(T_M^s\right)^2 \cdot \omega \\
    & \leqslant \int\left|u-u^B\right|^2 \cdot \omega.
    \end{aligned}
  \end{equation}
  The inequality in the above formula is due to $0 \leqslant T_M^s \leqslant 1$.
  Recall that the weight function $\omega = |x|^{2\mu}$, let
  \begin{equation}\label{eq:1111-0.5}
    y=B x, \quad d(y)=u(x)-u^B(x),
  \end{equation}
  then 
  \begin{equation}\label{eq:1111-1}
    \begin{aligned}
     \int\left|u-u^B\right|^2 \cdot \omega d x 
    &=  B^{-2 \mu-1} \int|d(y)|^2|y|^{2 \mu} d y \\
    &\leqslant  B^{-2 \mu-1} \int|d(y)|^2\bigl(1+y^2\bigr)^\mu d y.
    \end{aligned}
  \end{equation}
  Combining this with the definition of the $H^\mu$ norm, we have
  \begin{equation}\label{eq:1111-2}
    \int\left|u-u^B\right|^2 \cdot \omega d x \lesssim B^{-2 \mu-1}\bigl\|\mathcal{F}[d](\xi)\bigr\|_{H^\mu}^2.
  \end{equation}
  From \eqref{eq:0505-7} we have
  \begin{equation}\label{eq:0505-22}
    \mathcal{F}\bigl[u^B\bigr](k)=\mathcal{F}[u](k) \cdot T_B^f(k),
  \end{equation}
  then combining \eqref{eq:1111-0.5} yields
  \begin{equation}\label{eq:1111-3}
    \mathcal{F}[d](\xi)=B \cdot \mathcal{F}[u](B \xi) \cdot \left(1-T_1^f(\xi)\right).
  \end{equation}
  Putting \eqref{eq:1111-3} back into \eqref{eq:1111-2}, we have
  \begin{equation}\label{eq:1111-4}
    \int \left|u-u^B\right|^2 \cdot \omega d x \lesssim B^{-2 \mu+1} \left\| \mathcal{F}[u](B \xi)\left(1-T_1^f(\xi)\right)\right\|_{H^\mu}^2.
  \end{equation}
  Let $g(\xi) = \mathcal{F}[u](B\xi)$, we assert that the following conclusion holds:
  \begin{equation}\label{eq:1111-5}
    \left\|g(\xi)\left(1-T_1^f(\xi)\right)\right\|_{H^\mu(\mathbb{R})} \lesssim_\mu \bigl\|g(\xi)\bigr\|_{H^{\mu}\left(\mathbb{R}\backslash [-1,1]\right)}.
  \end{equation}
  By the definition of $T_B^f(\xi)$ \eqref{eq:0505-6} we have
  \begin{equation}\label{eq:0505-23}
    0 \leqslant T_1^f(\xi) \leqslant 1; \quad T_1^f(\xi) \equiv 1,\ \forall\,|\xi| \leqslant 1.
  \end{equation}
  Hence for $\mu \in \mathbb{N}$, we have
  \begin{equation}\label{eq:1111-6}
    \left\|g\left(1-T_1^f\right)\right\|_{H^\mu(\mathbb{R})}^2=\sum_{j=0}^\mu \int_{\mathbb{R} \backslash[-1,1]}\left|\partial^j\left(g\left(1-T_1^f\right)\right)\right|^2 d \xi.
  \end{equation}
  Combining
  \begin{equation}\label{eq:1111-7}
    \begin{aligned}
    \left|\partial^j\left(g\left(1-T_1^f\right)\right)\right|^2 & =\left|\sum_{r=0}^j\binom{j}{r}\left(\partial^r g\right) \cdot\left(\partial^{j-r}\left(1-T_1^f\right)\right)\right|^2 \\
    & \lesssim_\mu \sum_{r=0}^\mu\left|\partial^r g\right|^2
    \end{aligned}
  \end{equation}
  yields
  \begin{equation}\label{eq:1111-8}
    \left\|g\left(1-T_1^f\right)\right\|_{H^\mu(\mathbb{R})}^2 \lesssim_\mu \sum_{j=0}^\mu \int_{\mathbb{R} \backslash[-1,1]}\bigl|\partial^j g\bigr|^2 d \xi.
  \end{equation}
  That is to say, \eqref{eq:1111-5} holds for $\mu \in \mathbb{N}$.
  It then follows from the interpolation theory of Sobolev spaces that \eqref{eq:1111-5} holds for any $\mu \geqslant 0$.
  Substituting \eqref{eq:1111-5} into \eqref{eq:1111-4} then into \eqref{eq:0505-17}, we obtain
  \begin{equation}\label{eq:1111-9}
    E_f \lesssim_\mu B^{-\mu+\frac{1}{2}} \bigl\|\mathcal{F}[u](B \xi)\bigr\|_{H^{\mu}\left(\mathbb{R}\backslash [-1,1]\right)}.
  \end{equation}
  Next we consider the estimation of $E_H$ defined in \eqref{eq:0505-9}.
  Let
  \begin{equation}\label{eq:0505-33}
    g_{\xi, s}(x)=e^{-\frac{1}{2}(x-s)^2+i \xi x},
  \end{equation}
  then
  \begin{equation}\label{eq:0505-34}
    \begin{aligned}
    u_M^B & =u^B \cdot T_M^s \\
    & =\frac{1}{2 \pi} \int_{|\xi| \leqslant 2 B,|s| \leqslant 2 m} \mathcal{F}\bigl[u^B\bigr](\xi) g_{\xi, s}(x) d \xi d s.
    \end{aligned}
  \end{equation}
  From \ref{lem:transform} we know:
  let $z = \xi - is$, then the generalized Hermite expansion of $g_{\xi,s}$
  \begin{equation}\label{eq:0505-35}
    g_{\xi, s}(x)=\sum_n c_n \widehat{H}_n^{(\mu)}
  \end{equation}
  satisfy: when $n = 2m,\, m \in \mathbb{N}$,
  \begin{equation}\label{eq:0505-36}
    \begin{aligned}
    c_{2 m} & =\sqrt{\pi} \times(-1)^m \times \sqrt{\frac{\Gamma\left(m+\mu+\frac{1}{2}\right)}{m!}} \times \frac{1}{\Gamma\left(m+\frac{1}{2}\right)} \times\left(\frac{z}{2}\right)^{2 m} \\
    &{}\qquad \times e^{-\frac{z^2}{4}-\frac{1}{2} s^2} \times \phi\left(-\mu, m+\frac{1}{2} ; \frac{z^2}{4}\right),
    \end{aligned}
  \end{equation}
  when $n = 2m+1,\, m \in \mathbb{N}$,
  \begin{equation}\label{eq:0505-37}
    \begin{aligned}
    c_{2 m+1} & =\sqrt{\pi} \times(-1)^m \times \sqrt{\frac{\Gamma\left(m+\mu+\frac{3}{2}\right)}{m!}} \times \frac{1}{\Gamma\left(m+\frac{3}{2}\right)} \times i \times\left(\frac{z}{2}\right)^{2 m+1} \\
    &{}\qquad \times e^{-\frac{z^2}{4}-\frac{1}{2} s^2} \times \phi\left(-\mu, m+\frac{3}{2} ; \frac{z^2}{4}\right).
    \end{aligned}
  \end{equation}
  Here, $\phi\left(a,c;x\right)$ in \eqref{eq:0505-36} and \eqref{eq:0505-37} is the so-called
  Kummer's confluent hypergeometric function, which is defined as
  \begin{equation}\label{eq:0505-38}
    \phi(a, c ; x)=1+\frac{a}{c} \frac{x}{1!}+\frac{a(a+1)}{c(c+1)} \frac{x^2}{2!}+\ldots
  \end{equation}
  Since $z = \xi - is$, when $|\xi|,|s| \leqslant 2 B=2 M=\sqrt{\frac{N}{3}}$, we have
  \begin{displaymath}
    \left|\frac{z^2}{4}\right| \leqslant \frac{N}{6}.
  \end{displaymath}
  Then for $n > N,\, m=\left[\frac{n}{2}\right]$ we have
  \begin{equation}\label{eq:0505-39}
    \left|\frac{z^2}{4}\right| \leqslant \frac{m}{3}.
  \end{equation}
  Let $a = -\mu,\,c = m+\frac{1}{2}$ or $c = m+\frac{3}{2},\,x=\frac{z^2}{4}$,
  from \eqref{eq:0505-39} we know $|x / c| \leqslant \frac{1}{3}$, hence
  \begin{equation}\label{eq:0505-40}
    \begin{aligned}
    \left|\frac{a \cdots(a+k-1)}{c \cdots(c+k-1)} \frac{x^k}{k!}\right| 
    & =  \left|\frac{(-\mu) (-\mu+1) \cdots (-\mu+k-1)}{1 \cdot 2  \cdots  k}\right| \cdot\left|\frac{x}{c}\right| \cdot\left|\frac{x}{c+1}\right|  \ldots \left|\frac{x}{c+k-1}\right| \\
    & \leqslant  \frac{1}{3^k}.
    \end{aligned}
  \end{equation}
  From \eqref{eq:0505-40} we know that for $z=\xi-i s$ with $|\xi|,|s| \leqslant 2 B=2 M=\sqrt{\frac{N}{3}}$
  and $m=\left[\frac{n}{2}\right]$ with $n>N$, the following inequality holds:
  \begin{equation}\label{eq:0505-41}
    \left|\phi\left(-\mu, m+\frac{1}{2} ; \frac{z^2}{4}\right)\right|,\left|\phi\left(-\mu, m+\frac{3}{2} ; \frac{z^2}{4}\right)\right| \leqslant \sum_{k\geqslant 0} \frac{1}{3^k} = \frac{3}{2}.
  \end{equation}
  When $n>N$, combining \eqref{eq:0505-36}, \eqref{eq:0505-37} and
  \eqref{eq:0505-41} yields
  \begin{equation}\label{eq:0505-42}
    \left|c_n\right|^2 \lesssim e^{-\frac{1}{2} \xi^2-\frac{1}{2} s^2} \cdot \frac{\Gamma\left(\left[\frac{n+1}{2}\right]+\mu+\frac{1}{2}\right)}{\left[\frac{n}{2}\right]!} \cdot \frac{1}{\left\{\Gamma\left(\left[\frac{n+1}{2}\right]+\frac{1}{2}\right)\right\}^2} \cdot \left(\frac{|z|}{2}\right)^{2 n}.
  \end{equation}
  Further applying Stirling's formula yields
  \begin{equation}\label{eq:0505-43}
    \left|c_n\right|^2 \lesssim_\mu e^{-\frac{1}{2}\left(\xi^2+s^2\right)} \cdot n^\mu \cdot \frac{1}{n!} \cdot\left(\frac{\xi^2+s^2}{2}\right)^n.
  \end{equation}
  From \eqref{eq:0505-35} and \eqref{eq:0505-43} we know: when $|\xi|,|s| \leqslant 2 B=2 M=\sqrt{\frac{N}{3}}$,
  \begin{equation}\label{eq:0505-44}
    \begin{aligned}
    \left\|g_{\xi, s}-\widehat{\Pi}_N g_{\xi, s}\right\|_\omega^2 & =\sum_{n>N}\left|c_n\right|^2 \\
    & \lesssim_\mu \sum_{n>N} e^{-\frac{1}{2}\left(\xi^2+s^2\right)} \cdot n^\mu \cdot \frac{1}{n!} \cdot\left(\frac{\xi^2+s^2}{2}\right)^n \\
    & \lesssim_\mu \frac{\left(\frac{\xi^2+s^2}{2}\right)^{N+1}}{(N+1)!} \cdot N^\mu.
    \end{aligned}
  \end{equation}
  From \eqref{eq:0505-34} we have
  \begin{equation}\label{eq:0505-45}
    u_M^B-\widehat{\Pi}_N u_M^B=\frac{1}{2 \pi} \int_{|\xi| \leqslant 2 B,|s| \leqslant 2 M} \mathcal{F}\left[u^B\right](\xi)\left(g_{\xi, s}-\widehat{\Pi}_N g_{\xi, s}\right) d \xi d s.
  \end{equation}
  Applying the triangle inequality yields
  \begin{equation}\label{eq:0505-46}
    \left\|u_M^B-\widehat{\Pi}_N u_M^B\right\|_\omega \leqslant \frac{1}{2 \pi} \int_{|\xi| \leqslant 2 B,|s| \leqslant 2 M}\left|\mathcal{F} \left[ u^B \right](\xi)\right|\left\|g_{\xi, s}-\widehat{\Pi}_N g_{\xi,s}\right\|_\omega.
  \end{equation}
  Further applying the Cauchy-Schwarz inequality yields
  \begin{equation}\label{eq:0505-47}
    \begin{aligned}
    \left\|u_M^B-\widehat{\Pi}_N u_M^B\right\|_\omega^2 & \lesssim\left(\int_{|\xi| \leqslant 2 B,|s| \leqslant 2 M}\left|\mathcal{F}\left[u^B\right](\xi)\right|^2 d \xi d s\right) 
    \times\left(\int\left\|g_{\xi, s}-\widehat{\Pi}_N g_{\xi, s}\right\|_\omega^2 d \xi d s\right) \\
    & \lesssim 4 M\|u\|^2 \int_{|\xi| \leqslant 2 B,|s| \leqslant 2 M}\left\|g_{\xi, s}-\widehat{\Pi}_N g_{\xi, s}\right\|_\omega^2.
    \end{aligned}
  \end{equation}
  From \eqref{eq:0505-44} we have
  \begin{equation}\label{eq:0505-48}
    \begin{aligned}
    & \int_{|\xi| \leqslant 2 B,|s| \leqslant 2 M}\left\|g_{\xi, s}-\widehat{\Pi}_N g_{\xi, s}\right\|_\omega^2 d \xi d s \\
    &\lesssim_\mu \int_{|\xi| \leqslant 2 B,|s| \leqslant 2 M} \frac{1}{(N+1)!} \cdot\left(\frac{\xi^2+s^2}{2}\right)^{N+1} \cdot N^\mu d \xi d s \\
    & \leqslant \frac{N^\mu}{(N+1)!} \cdot \int_{\left(\xi^2+s^2\right) \leqslant 4\left(B^2+M^2\right)}\left(\frac{\xi^2+s^2}{2}\right)^{N+1} d \xi d s \\
    &= \frac{N^\mu}{(N+1)!} \cdot 2 \pi \cdot 2^{-(N+1)} \cdot \frac{N^{N+2}}{2 N+4} \cdot\left(\frac{2}{3}\right)^{N+2} \\
    &\lesssim_\mu N^{\mu-\frac{1}{2}} e^{-\frac{N}{12}}.
    \end{aligned}
  \end{equation}
  Putting \eqref{eq:0505-48} back into \eqref{eq:0505-47} we know
  $E_H$ defined in \eqref{eq:0505-9} satisfies:
  \begin{equation}\label{eq:0505-49}
    \begin{aligned}
    E_H&=\left\|u_M^B-\widehat{\Pi}_N u_M^B\right\|_\omega \\
    &\lesssim_\mu e^{-\frac{N}{24}} M^\mu \|u\|.
    \end{aligned}
  \end{equation}
  Combined the estimation of $E_s$ in \eqref{eq:0505-16}, the estimation
  of $E_f$ in \eqref{eq:1111-9}, the estimation of $E_H$ in \eqref{eq:0505-49}
  with \eqref{eq:0505-9}, we prove \eqref{eq:bigthm} under the condition $\beta = 1$.
  As discussed before, by further applying \eqref{eq:0505-2}, we can prove the general
  case when $\beta \neq 1$, this completes our proof.
\end{proof}

\end{document}
